%%%%%%%%%%%%%%%%%%%%%%%%%%%%%%%%%%%%%%%%%%%%%%%%%%%%%%%%%%%%%%%%%%%%%%%%%%%%%%%%
%2345678901234567890123456789012345678901234567890123456789012345678901234567890
%        1         2         3         4         5         6         7         8

\documentclass[letterpaper, 10 pt, conference, dvipdfmx]{ieeeconf}  % Comment this line out if you need a4paper

%\documentclass[a4paper, 10pt, conference]{ieeeconf}      % Use this line for a4 paper

\IEEEoverridecommandlockouts                              % This command is only needed if 
                                                          % you want to use the \thanks command

\overrideIEEEmargins                                      % Needed to meet printer requirements.

%In case you encounter the following error:
%Error 1010 The PDF file may be corrupt (unable to open PDF file) OR
%Error 1000 An error occurred while parsing a contents stream. Unable to analyze the PDF file.
%This is a known problem with pdfLaTeX conversion filter. The file cannot be opened with acrobat reader
%Please use one of the alternatives below to circumvent this error by uncommenting one or the other
%\pdfobjcompresslevel=0
%\pdfminorversion=4

% See the \addtolength command later in the file to balance the column lengths
% on the last page of the document

% The following packages can be found on http:\\www.ctan.org
%\usepackage{graphics} % for pdf, bitmapped graphics files
%\usepackage{epsfig} % for postscript graphics files
%\usepackage{mathptmx} % assumes new font selection scheme installed
%\usepackage{times} % assumes new font selection scheme installed
%\usepackage{amsmath} % assumes amsmath package installed
%\usepackage{amssymb}  % assumes amsmath package installed
\usepackage{graphicx}
%\usepackage{multirow}
%\usepackage{multicol}
%\usepackage{comment}
%\usepackage{breqn}
%\usepackage{bm}
%\usepackage{color}
%\usepackage{diagbox}
\usepackage{url}

\usepackage{amsmath, amssymb}
%\usepackage[top=20.1truemm,bottom=15.2truemm,left=16.9truemm,right=16.9truemm]{geometry}
% \usepackage[top=19.1truemm,bottom=19.1truemm,left=19.1truemm,right=19.1truemm]{geometry}

\title{\LARGE \bf
Deep Predictive Learning : Motion Learning Concept\\
inspired by Cognitive Robotics
}
\author{ Kanata Suzuki$^{1,2,\dagger}$, Hiroshi Ito$^{1,3,\dagger}$, Tatsuro Yamada$^{1,4,\dagger}$, Kei Kase$^{1,\dagger}$ and Tetsuya Ogata$^{1,\dagger}$
\thanks{$^{1}$All authors are with Waseda University, Tokyo, 169-8555, Japan
        {\tt\small ogata@waseda.jp}}% <-this % stops a space
\thanks{$^{2}$Kanata Suzuki is with Fujitsu Limited, Tokyo, 105-7123, Japan}
\thanks{$^{3}$Hiroshi Ito is with Hitachi, Ltd., Ibaraki, 312-0034, Japan}
\thanks{$^{4}$Tatsuro Yamada is with Panasonic Corporation, Osaka, 571-8501, Japan}
\thanks{${\dagger}$ All authors have contributed equally to this work.}
}

 % Improvement of interpretability and disturbance robustness of deep predictive learning by modality attention

\begin{document}


\maketitle
\thispagestyle{empty}
\pagestyle{empty}

\begin{abstract}
A deep learning-based approach can generalize model performance while reducing feature design costs by learning end-to-end environment recognition and motion generation. However, the process incurs huge training data collection costs and time and human resources for trial-and-error when involving physical contact with robots. Therefore, we propose ``deep predictive learning,'' a motion learning concept that assumes imperfections in the predictive model and minimizes the prediction error with the real-world situation. Deep predictive learning is inspired by the ``free energy principle and predictive coding theory,'' which explains how living organisms behave to minimize the prediction error between the real world and the brain. Robots predict near-future situations based on sensorimotor information and generate motions that minimize the gap with reality. The robot can flexibly perform tasks in unlearned situations by adjusting its motion in real-time while considering the gap between learning and reality. This paper describes the concept of deep predictive learning, its implementation, and examples of its application to real robots. The code and document are available at \url{https://ogata-lab.github.io/eipl-docs}
\end{abstract}

%%%%%%%%%%%%%%%%%%%%%%%%%%%%%%%%%%%%%%%%%%%%%%%%%%%%%%%%%%%%%%%%%%%%%%%%%%%%%%%%
\section{Introduction}
The ``Moravec Paradox'' \cite{moravec1988mind} is one of the greatest challenges in artificial intelligence technology. This paradox refers to the contradiction that tasks that can be performed semi-consciously by children, compared to advanced logical reasoning performed consciously by adults (e.g., puzzles and chess), are difficult for modern artificial intelligence and robots. In recent years, with the advent of deep learning, significant technological progress has been made in locomotion, such as walking and running, in the 2000s, and cognition, such as vision and hearing, in the 2010s. However, despite today's state-of-the-art technologies, it is challenging for a robot to perform various tasks with its manipulator that requires interaction between the robot and its environment. This paper describes the concept and implementation of deep predictive learning, a method for applying deep learning techniques to robotics. Chapter 3 describes robotics applications using deep learning and their challenges. Chapter 4 briefly explains the background of deep predictive learning, including predictive coding, dynamic cognitive systems, and subsumption architecture. Chapter 5 describes a concrete implementation of deep predictive learning, and Chapter 6 describes an actual robotics application case study. Finally, Chapter 7 provides a summary and future perspectives.

\section{Robot Applications of Deep Learning}
The term ``deep learning'' was first mentioned at a top robotics conference in 2017 at IEEE/RSJ IROS 2017, much later than other fields.
Six years later, the application of deep learning to robotics remains limited compared to the fields of image processing and natural language processing. It is primarily utilized for image recognition, speech recognition, and deep reinforcement learning, and not much for motion learning.

\subsection{Recognition}
Most early deep learning applications focused on object recognition (robot vision) and speech recognition (dialogs). Traditional robot intelligence design typically involves the following steps: (1) recognize the environment and objects, (2) map the recognition results to a 2D map or 3D CAD model manually pre-deigned by a human, (3) determine the relative pose and location of the target object, and (4) plan its movements according to the situation. Deep learning has greatly improved the first step, the recognition capability. By feeding manually pre-designed labels such as suitable grasp pose to the deep learning framework, it can recognize and generalize to various objects that were known to be difficult to model. However, the burden of human labeling is significant, and learning from images alone does not guarantee the successful grasping of objects. This is because grasping requires not only vision but also needs to consider object attributes and physical features of the robot. Such considerations include the gripper's grasping force, friction conditions, and object deformation.

\subsection{Deep Reinforcement Learning}
Many researchers are currently working on applying deep reinforcement learning to robots. After its success in learning games, its application to robots took place naturally. 
In reinforcement learning, the robot is trained to predict expected reward from a sensorimotor input such as images, and from its behavioral motion output. The main advantage of this approach is that it considerably reduces the process of designing models of the environment, which traditionally was manually scripted. However, reinforcement learning frameworks typically require large training data to obtain an ``optimal model.''

The Sim2Real approach is one of the ways to negate the problem of collecting large training data in the real world. 
With the advancement of simulators, it can generate visually realistic world, but modeling physics and contact conditions remains challenging.
This issue is currently discussed as a single discipline called ``surface science.''


\section{Deep Predictive Learning}

\subsection{Concept of Predictive Coding}
The machine learning framework involves learning an optimal model with a large amount of data and ensuring that the model works appropriately on new data, i.e., has maximum generality. This concept is the same for object recognition and reinforcement learning. However, obtaining an optimal model is challenging, considering the impracticality of predicting which state will give maximum generalizability. Especially robots can be assumed to be constantly exposed to unexpected situations to operate in the real world. Therefore, adopting an approach that considers learning as imperfect and deals with this error is necessary. A possible framework is Friston's free energy principle~\cite{friston2012perception}, given by 

\begin{equation}
    FreeEnergy = -log(x|z) + D_{KL}[q(z|x)||p(z)] \label{eq_free}
\end{equation}

This equation corresponds to the negative evidence lower bound (ELBO) in general machine learning models, which is not new. However, based on the ``interpretation'' of this equation and its multiple parameters, we can consider designing an embodiment system.

The brain is a predictor and has as a prior belief $p(z)$ about $z$ that causes the sensation $x$. Here, perception $z$ includes the hidden states (causes) of exteroceptive sensations such as vision and touch and the hidden states of proprioceptive sensations (e.g., joints) and endoceptive sensations related to motor and visceral sensations. The brain predicts sensation $x$ based on perception $z$ and receives the current proprioceptive sensory input $x$ to obtain the posterior belief $q(z|x)$. A new prediction for the next sensation $x$ is generated based on the updated posterior belief. The prediction generation and belief updating are performed continuously. The predicted sensation $x$ includes a prediction of accuracy (the inverse of the variance of the sensation $vx$). The first term in equation \ref{eq_free} refers to the prediction error (Accuracy) weighted by this accuracy, and the second term refers to the effect of the complexity of the prior belief (model) structure (Complexity).

Strategies to reduce prediction error are called ``predictive coding,'' and four main strategies are possible.

First, learn by training the entire model that generates the sensory predictions and changing the internal parameter $w$ according to the error. However, this approach is costly and challenging to implement in time-constrained situations. Furthermore, there may be cases where frequent model modifications are inappropriate.

Second, modify the posterior belief $q(z)$. For example, when humans perceive the environment, they may comfortably interpret the real world according to the prior images in their brains and are unaware of the error. This method uses predictive biases about the world to adapt to changes in the real world by creating a virtual perception that differs from reality.

Third, use the body to influence the external real world by changing sensory $x$ to reduce prediction error. We achieve this by manipulating the environment or changing one's perspectives. This strategy is essential for autonomous agents, such as robots, with bodies in the real world.

Fourth, reduce the apparent error by setting a smaller prediction accuracy of the sense $x$, $1/vx$, and multiplying it by the ``weight'' of the prediction error.
This is related to the concept of ``attention'' discussed below.




\subsection{Comparison with Reinforcement Learning}
This section discusses the main differences between elementary reinforcement learning and predictive coding. 

One of the distinctive characteristics of reinforcement learning is that the learning goal can be set freely and abstractly (e.g., scalar values).
In contrast, the main goal in predictive coding is to minimize the prediction error in the model's connection between the body and the world.
Therefore, in principle, it is difficult to set other learning goals in predictive coding unless special efforts are made.

In model-based reinforcement learning, the ``world model'' is learned through repeated interactions between actions and states.
In \cite{ha2018world}, actions and states are learned as a set of abstracted latent spaces; the idea is similar to predictive coding in that the world is the target of the prediction.
However, the world model in reinforcement learning searches for the action policy that can achieve the maximum reward in the model by trial and error after sufficient learning has progressed.
Thus, the optimal behavior based on the learned world model can be acquired in reinforcement learning.

In contrast, the predictive coding model assumes an open world and that learning the model of the environment is continuous.
With insufficient learning, the model generates actions that are expected to have minor future prediction errors (expected free energy).
Then, by repeating the experience of reaching the target state, the sequence of actions becomes ``habituated'' to the prior belief $p(x)$.
The generated series of behavioral sequences (virtual expectations) is based on trial and error in this habituation process.
However, the learned motions need not be optimal for any given objective measure.

For example, right-handed people pick objects in front of them using their right hand, even if the object is slightly to the left.
They do not force themselves to pick it up with their left hand. This behavior is habitual, i.e., they choose an action with high predictive accuracy.
In the next section, we discuss implementing a mechanism that partially reflects this idea of predictive coding in a robot.


\subsection{Physicality Perspective}
An interesting consideration can be drawn from discussing physicality as the norm, as exemplified by behavior-based robotics~\cite{brooks1991intelligence}.
Here, it is assumed that intelligent physical behavior emerges from the coupling (mechanical attraction) of the ``environment-body-nervous system.''
In particular, it emphasizes the interaction between the body and the environment.
For example, in the passive walking mechanism, ``walking'' emerges from the interaction between a cleverly designed bipedal (or multipedal) articulation model and a slope.
We also focus on the various ``functions'' that emerge from the interaction between the skillfully designed body and the environment, including the body's flexibility, mechanism, and sensor placement.
In other words, the body itself implicitly retains the motion and actions.
At this point, the nervous system is considered to perform only minimal functions, in an example of passive walking, the oscillator system of the body is enough to execute the action.


\subsection{Dynamic Cognitive System Perspective}\label{sec_cognitive}
Given the ``environment-body-nervous system'', the nervous system can have multiple attractor states.
An example is the discussion of dynamic cognitive systems based on the dynamical systems perspective proposed by Tani ~\cite{tani1996model}.
Here, the brain is assumed to be an input-output and predictive system for the world (Figure \ref{fig_attractor}).
In interaction with the body, the brain sometimes forms a specific attractor structure with the world, and when the discrepancy grows, the connection is broken.
The brain then transitions to a new attractor, considered a phase transition, sometimes through a chaotic state.
As this process repeats, cognition never converges but transitions dynamically.
Tani et al. realized this process using a recurrent neural network (RNN) and a real robot system and have been continuously doing exciting research to understand the cognitive processes of the brain.
Although this study was before the introduction of predictive coding and the free energy principle, the concepts have much in common.


%-------------------- Fig --------------------
\setlength\textfloatsep{5pt}
\begin{figure}[t]
\centering
\includegraphics[width=\columnwidth]{attractor.eps}
\caption{Framework of a dynamic cognitive system}
\label{fig_attractor}
\hspace{10mm}
\end{figure}    


\subsection{Examples of Deep Predictive Learning} \label{sec_attention}
With the above ideological background, we believed that the ``neural system'' part of the ``environment-body-neural system''coupling could be configured in a deep learning model to couple the robot with the world to perform higher dimensional and more complex tasks. The first model was published in 2014 \cite{noda2014multimodal}, and since then, this area of research has continued to expand.

Deep predictive learning combines multiple deep learning models to build a sensory and motion prediction system. For example, convolutional autoencoder and RNNs are combined to build a sensory-motor prediction system coupled with the environment and body. The following are some critical structural points of the deep predictive learning model. The sensation and motion are integrated into a single predictive system. In a typical engineering application, mapping is learned with sensation as the input and motion as the output. However, in the deep predictive learning framework, sensation and motion form a single perceptual system that is both the input and the object of prediction. The framework is based on learning this perceptual continuum as an experience. Using RNNs with different time scales is also essential, as shown in figure \ref{fig_mtrnn}. A hierarchical approach is used where the lower layer, closer to the senses, predicts the early (short time constant) dynamics, while the upper layer predicts the late (long time constant) dynamics. Methods such as multiple-time scale RNN (MTRNN)\cite{yamashita2008emergence} and long short-term memory (LSTM)\cite{Hochreiter1997nc} are used to implement this, and the interactions help minimize the overall prediction error. Specifically, the prediction error of the senses is absorbed by the state changes in the lower layers. Here, the upper layer generates predictions based on long-term expectations (intentions) and provides constraints to the lower layer to maintain this state. However, if the lower layer cannot absorb the prediction errors of the senses, the errors affect the upper layer, changing the state of the upper layer (change in intention). Thus, the structure in which the upper and lower layers interact allows generative adaptive behavior while maintaining consistency. From a hierarchical perspective, the predictive coding inspired RNN (PV-RNN) by Tani et al. is the latest predictive coding model that allows updating the prior belief $p(z)$ at each level, including the upper and lower layers)\cite{ahmadi2019novel}. This model deals aggressively with prediction error minimization using real-time error backpropagation (called regression) and has high theoretical consistency. Although some problems, such as computational cost, exist, applying this model to real robots with multidimensional degrees of freedom can be beneficial. In addition, a recent study investigated the equivalence of hierarchical neural networks with recursive coupling and predictive coding\cite{isomura2022canonical}.




%-------------------- Fig --------------------
\setlength\textfloatsep{5pt}
\begin{figure}[t]
\centering
\includegraphics[width=60mm]{S-MTRNN.eps}
\caption{Hierarchical structure of recursively coupled predictors with different time scales}
\label{fig_mtrnn}
\hspace{10mm}
\end{figure}


\newpage

\setlength\textfloatsep{5pt}
\begin{figure}[t]
\centering
\includegraphics[width=\columnwidth]{attention_bag.eps}
\caption{Bag opening task using an attention mechanism}
\label{fig_attention_bag}
\hspace{10mm}
\end{figure}   


A more detailed structure and implementation methods of deep predictive learning are described in Section \ref{sec_implementation}. The following gives an overview of deep predictive learning using a concrete example. Figure \ref{fig_attention_bag} shows a robot~\footnote{https://youtu.be/SC0BfTlnYmw} that can open the zipper of a cloth bag, which the authors presented in ~\cite{ichiwara2022contact, ito2022integrated}. Here, the cloth bag, the robot arm (with a camera and tactile sensors), and the deep learning model are retracted so that the perceptual prediction (posterior belief) $q(z)$ in the model is modified in real-time (i.e., the predicted image, tactile state, and motion are modified). The input sensation $x$ is also modified by motion modification, and the real-time model prediction error is reduced. In conventional control technology, the trajectory of the robot arm is given in advance as a target for real-time error correction. However, in the ``zipper opening task'' discussed in this paper, generating an appropriate trajectory in real-time according to the state of the bag and robot hand is necessary. Furthermore, as this task starts with grasping a randomly placed zipper, the grasping position and posture constantly change, increasing the task's difficulty level. In this study, direct teaching is used to train the robot to open the zipper multiple times to give the robot a real experience with the task we want to convert into a habit. The number of times the movement is taught is 36, which is extremely small compared to common learning methods.


\begin{description}
\item[Trial and error]\mbox{}\\
Owing to online motion generation by the real robot using the trained model, granular trajectory modification and motion switching were based on a trial-and-error process.
Specifically, this means that the robot generates a motion based on the prediction (prior belief) $p(x)$. The motion is then modified in real-time (i.e., a transition to another motion attractor) using a retraction based on the prediction error.
The trial and error is not random but a modification of the predicted motion based on past experiences.


\item[Bottom-up and top-down attention]\mbox{}\\
By introducing an ``attention mechanism,'' the model becomes robust to sudden changes in lighting conditions, placement of multiple objects as noise, and movement of the target object, among others. Specifically, based on verbal attention instructions, we confirm the phenomenon of ``disappearance of backgrounds other than the attended object'' in the generated predictive images (posterior beliefs) ~\cite{ito2022integrated}. It is important to note that attention here is created by ``prediction.'' Normally, attention in image processing is generated from the ``current input image'' in a bottom-up manner. It involves selecting features such as luminance, color, and orientation using saliency-based feature selection of features or by generating a mask, key, query, and value from the image to maximize the detection reward. The specific implementation of top-down attention prediction is under debate. In the zipper-opening task, several points of attention are predicted in the image. In other words, we limit the image regions that can be attended to and ignore the rest. The loss function is a prediction error that includes all visual, tactile, and motor components; thus, the points of attention are self-organized into predictable parts that change (can be manipulated) depending on one's actions.
\end{description}


\setlength\textfloatsep{5pt}
\begin{figure*}[t]
\centering
\includegraphics[width=17cm]{1_overview.eps}
\caption{
Example of a deep predictive learning framework consisting of a feature extractor and time-series predictor.
}
\label{fig_overview}
\end{figure*}


\section{Implementation}
\label{sec_implementation}
This section describes the detailed structure and implementation methods of deep prediction learning. Figure \ref{fig_overview} shows the overall structure of deep prediction learning with three phases: training data collection, learning, and motion generation. The following subsections describe the outline and features of each phase.

\subsection{Data collection}
In deep predictive learning, motion generation models are acquired in a data-driven manner using the robot's sensorimotor information (time-series data comprising motion and sensor information) as training data. Therefore, the training data must contain information about the interaction between the robot's body and the environment. In addition, since the quantity and quality of the training data influence model performance, high-quality demonstration data must be collected efficiently. In Phase 1, the training data is collected by performing the desired motions on the robot and recording motion information, such as joint angles, and sensor information, such as camera images, at a constant sampling rate. Typical teaching methods for robot motion include program description, direct teaching, and teleoperation. Although describing the robot's motion in advance using a robot programming language is simple, it may not be feasible due to the complexity of the description when a robot needs to perform a long-term motion. In contrast, training data without precise modeling and parameter adjustment can be obtained by teaching movements through human manipulation of the robot. Among them, remotely controlling the manipulator from the actual robot's perspective (Wizard of Oz \cite{Koma2016ral}) is desirable because it can intuitively teach the robot human manipulation skills for a task. The operator interacts naturally with the environment by controlling the robot as if he were controlling his body. In addition, since the operator makes decisions about actions based on information obtained from sensor information, the acquired teaching dataset is expected to contain information necessary for motion learning and be effective in acquiring data for model training.


\subsection{Integrated learning of sensorimotor information}
In deep predictive learning, the learning target is the time-series relationships between sensorimotor information in a system in which the environment and the body interact dynamically. The training data are not labeled with correct answers, and the model is trained to predict the robot state in the next step using the current robot state as input. This autoregressive learning eliminates the need for the detailed design of a physical model of the environment, as required in conventional robotics. In addition, the model can be represented as a dynamic system integrating environment recognition and motion generation functions across multiple modalities.

The model consists of feature extraction and time series learning parts to learn the robot's sensorimotor information. The feature extraction part extracts features from sensor information values acquired by the robot, and the time series learning part learns sensorimotor information that integrates the extracted features and robot motion information ( e.g., joint angles and torque). Although each part is connected end-to-end or learned independently, the roles of each part are explicitly separated in this chapter.


\subsubsection{Feature extractor}
Because the inputs the robot acquires using sensors, such as camera and tactile sensors, from the real environment are often high-dimensional, feature extraction is critical. To learn sensorimotor information appropriately, the extracted sensor features must reflect the relationship between the robot body and manipulated object. This is because it affects the prediction performance of the motion information in the time series learning section at a later stage, e.g., the approach to the manipulated object in the manipulation task. Therefore, to obtain detailed information, the feature extractor should be able to handle high-dimensional sensor inputs as they are. 

For example, consider raw image data containing unnecessary information for the robot task, such as background images, that become noise for model training. Most of our models use a convolutional autoencoder (CAE) structure as a feature extractor, which performs unsupervised learning so that the input and output are identity maps. In CAE, low-dimensional features are acquired in the middle layer by performing self-supervised learning to provide output values equal to the input values. The encoded feature vector represents the state of the input image and can be used to provide high-dimensional input information using fewer dimensions, making it robust to unwanted noise. Moreover, the above feature vectors can be further constrained to useful information for task learning by introducing an attention mechanism\cite{ichiwara2022contact, ito2022integrated}.


\subsubsection{Time-series predictor}
We use RNNs with recursive inputs and outputs in the middle layer as time series predictors. RNN models time dependency and functions as a learning-based controller that takes the current state as input and outputs the next state, including the motion trajectory. RNN has feature vectors obtained by the sensor feature extraction and robot motion information (such as joint angles, torque values, and gripper values) as input and output at each time and is trained to reduce the error between the output and input at the next time. A robust model can be constructed by adjusting the ratio of how much of the model output from the previous time is used as input. Moreover, the prediction of sensor values in a subtask manner contributes to the model's prediction performance. 

In addition, generating composite motions by representing long-term tasks as a combination of dynamics embedded in the learner is possible. RNNs can retain different dynamics obtained from the learning process of time-series data as unique internal states. During motion generation, the RNN predicts the next state based on the internal and current sensorimotor information, but the prediction error increases in the unlearned state. When the prediction error increases, the system switches from the learned dynamics to minimize the prediction error; this property can be used to switch modules. In addition, by explicitly designing an intermediate layer with multiple timescales, short- and long-term subtasks can be represented \cite{Hochreiter1997nc, Yamashita2008plos}. Thus, various task learning is possible using top-down and bottom-up hierarchical structures across multiple time scales and an engineering-like design of internal states. More detailed examples are described in section \ref{sec_case_studies}.


\subsection{Motion generation}
When performing the task, RNN performs three processes sequentially: (1) Acquire sensor information from the robot, (2) predict the next state based on the sensor information, and (3) send control commands to the robot based on the predicted values. By performing the forward computation of the model at each step, RNN predicts the robot's state for the next time-step based on the context information and inputs it holds internally. The RNN output is then used as the target state to control each joint. By repeating the above process online, the RNN predicts the sensorimotor motion of the robot while sequentially changing the state of each neuron in the context layer. Based on the results of this prediction and the prediction error in the real environment, the robot generates motions that dynamically respond to the input.

Another advantage of using deep learning models for motion generation is the online motion generation speed. The proposed framework comprises lightweight models, and the computational time and cost required for motion generation are low. Moreover, implementing each of the previous functions as components makes it possible to easily reuse the implemented system when tasks or robot hardware changes, or devices are added ~\cite{Kanamura2021sensors}.

% Learning Skills Policy
% Robot Applications
\section{Robot Applications}
\label{sec_case_studies}

While task accuracy has improved~\cite{Levine2017ijrr, Andrychowicz2020ijrr}, based on large-scale data and large numbers of trials, the dynamics of real-world environments are long-term and complex, requiring a wide range of task motions to be realized in a limited amount of motion data.
Because several practical tasks are essentially long-term tasks with multiple steps, an approach that divides and combines them into short-term operations is effective~\cite{Sacerdoti1974ai,Wolpert1998,Schaal2003isrr}.
The ``combined generation of robot motions'' improve the efficiency of training data collection and the versatility of the motion generation model.
\begin{itemize}
    \item Diversion of subtasks included in the learning task to other tasks (module reuse)
    \item Combination of subtasks to execute new tasks (extension of functions)
\end{itemize}

With the above background, this section discusses factors required to generate complex robot motions and presents examples of applications using real robots with deep predictive learning.

First, a simple approach to generating complex robotic motions is to coordinate each modularized function. In general, a robotic system consists of a set of multiple modules, and their coordination determines a set of outputs. In robot operation, it can be assumed that by combining or switching models with different functions in parallel, the overall task function can be extended without increasing the computational load of each model. Another advantage of the distributed system configuration is the ability to continuously add functions by modifying or adding modules.

Next is the local composite generation, which focuses on a single module. As the task is further subdivided, having the ability to perform multiple subtasks in each module is desirable. The dynamics must be represented in the module in a combinable format for composite multiple robot motions. By selecting the appropriate dynamics within the model in response to inputs to the model, it is possible to perform long-term task motion in response to environmental changes.

Finally, we ground each segmented subtask with its linguistic representation. By learning that each task execution is semantically grounded in a given linguistic instruction, the model can acquire internal dynamics as a representation, with the semantic and syntactic structure of the language system as a constraint. The resulting model can be deployed to use cases in which it sequentially generates motions while communicating with collaborators. In addition, linking one-time motion generation with discrete symbols having abstracted meanings in the real world with many disturbances and high situational dependence will bring benefits such as improved learning efficiency and robustness of motion generation and generalizability to operate in novel situations. Based on the above ideas, we aimed to inductively represent the relationship ~\cite{Harned1990pd} between motions and symbols in the sensorimotor information inside the model. In the following, we describe examples of robot applications for each element.


%%%%%%%%%%%%%%%%%%%%%%%%%%%%%%%%%%%%%%%%%%%%%%%%%%%%%%%%%%%%%%%%%%%%%%%%%%%%%%%%%%%%%%%%%%%%%%%%%%%%%%%%%%%%%%%%%%%%%%%%%%%%%

\subsection{Multiple model switching based on prediction error}
One method of switching subtask modules mentioned in element (i) is based on the prediction error of the learner. The model can hold various dynamics from the time series data learning process as a unique internal state. During motion generation, the next state is predicted based on the internal state and current sensorimotor information, but the prediction error tends to increase in unlearned states. Therefore, when the prediction error increases during motion generation, the system switches from the learned dynamics to minimize the prediction error; this property can be used for module switching. If a large prediction error cannot be handled in real-time, state transition that pulls in individual RNNs as described in Section \ref{sec_cognitive}, cannot be used alone. In such cases, it is necessary to coordinate multiple modules according to the prediction error size and adjust the overall sensory prediction and motion generation.

In \cite{Ito2022sr}, we developed a method to reduce prediction error in more complex tasks by introducing multiple prediction models and a novel switching mechanism between them according to prediction error (Figure \ref{fig_doorrobot}). The robot trains the prediction model for each motion, such as ``opening a door'' or ``passing through,'' and combines them to perform a series of tasks. Each prediction model uses information from sensors to generate a prediction image and compares it with the visual image of the actual robot (real image) to calculate in real time the change over time of an index (confidence), indicating how accurately the robot can work with the real state. In addition, the robot autonomously selects the prediction model with the highest confidence level to perform motions appropriate to the situation. Because the robot performs real-time calculations, it can handle complex tasks that cannot be handled by a single prediction model without having to precisely design the timing of motion switching and motion sequence. The verification results show that the robot can flexibly generate motions for unlearned door positions and changes in the color and pattern of door panels due to the retraction motion of the prediction models. In the task, ``opening a door,'' when a sudden change of situation occurred, the robot was moved to another location while in the process of opening a door. Normally, the robot would perform the motions of approaching the door, opening the door, and passing through the door in the same order. However, when a disturbance occurs, the robot immediately recognizes that the situation has changed and generates the motion of approaching the door again, eventually allowing the robot to pass through the door. Another task that changes the motion sequence is when a person opens the door that the robot is trying to open. The robot can perform a seamless operation by stopping the door-opening operation and switching to the door-passing operation when the door suddenly opens.


%-------------------- Fig --------------------
\setlength\textfloatsep{5pt}
\begin{figure}[t]
\centering
\includegraphics[width=73mm]{DoorRobotv2.eps}
\caption{Real-time motion switching of multiple prediction models based on prediction error}
\label{fig_doorrobot}
\hspace{10mm}
\end{figure}


In contrast, in \cite{Suzuki2021ral}, using multiple control modules with different characteristics, the authors addressed robot states (undefined motions) that cannot be predicted by learning-based controllers, including a prediction model and are not present in the training data. In the proposed method, a model-based controller compensates for the RNN output and switches it appropriately according to the error between the predicted and actual trajectories. Here, we assume that the internal state of the RNN contains the most recent information necessary to perform the task. Therefore, we train an additional task that directly predicts the past trajectory from the internal state. By adding a neuron to the internal state that predicts the additional task, we can add an error recovery function without affecting the performance of the main task learned by the RNN. Our experiments confirmed that error recovery is possible for different perturbations (contact and environmental changes) of cooperative operators by combining it with the dynamic design method of RNN internal states, which is described in the next section.

In general robot systems, much development and adjustment work is required because responses to unexpected situations must be programmed separately in addition to the expected sequence of operations. In addition, when an unexpected disturbance occurs in the working environment, the robot operates slowly because it has time to re-recognize and re-plan the situation. In \cite{Ito2022sr,Suzuki2021ral}, multiple controllers can be switched to cope with sudden changes in the situation, and operation procedures and approaches, such as the mixture-of-expert, based on these prediction errors are effective for autonomous task execution in response to environmental changes.




%%%%%%%%%%%%%%%%%%%%%%%%%%%%%%%%%%%%%%%%%%%%%%%%%%%%%%%%%%%%%%%%%%%%%%%%%%%%%%%%%%%%%%%%%%%%%%%%%%%%%%%%%%%%%%%%%%%%%%%%%%%%%
\subsection{Designing the dynamics within the model}
Next, we discuss the method of ``designing the dynamics within the model,'' corresponding to local composite generation within a single module as described in local composite generation. The methods described in the previous section focused on switching between modules, but for scalability, a mechanism that allows each module to switch its motion is desirable. We implement this by explicitly embedding sensorimotor dynamics in the learner on a subtask-by-subtask basis.

Although combining multiple subtask modules using an attractor structure is possible, the intended attractor structure must be embedded for switching. In \cite{Kase2018icra}, an attractor structure was embedded in the RNN using a loss function for the internal state of the RNN to combine multiple tasks. By constraining the attractor to have a common state that can be switched between the attractors of each subtask, an attractor pull-in motion was generated that reduced the prediction error in the designed switching state. In our experiments, we achieved motion generation in arbitrary subtask order by combining multiple subtasks for the put-in-box task\footnote{https://youtu.be/QwogmDqOxfE}. Furthermore, in \cite{Kase2019iros}, the loss function of the internal state of the RNN was extended to design an arbitrary attractor structure, thereby increasing the degree of freedom of the motion-switching posture.

In \cite{Suzuki2018ral}, a hierarchical structure of motion instruction input and RNN was introduced to learn subtask switching during a long-term task. By inputting a motion instruction vector that combines the collaborator's intention with the designed fixed-point attractor, it is possible to handle motion switching in ambiguous situations that cannot be judged based on sensor information alone. In the experiment, we verified a multi-motion folding task and executed the motion according to the motion instructions. In addition, we confirmed that the RNN's internal state has an attractor structure with motion instructions embedded in the upper long-term time-scale layer and motion in the short-term time-scale layer. In addition, it appropriately combines task context, text language, and image input.

While these methods of combining and switching motions using attractor structures require the prior design of subtasks, they can somewhat incorporate the designer's intentions. More flexible robot motion switching is expected when combined with grounding using the language described below.



%-------------------- Fig --------------------
\setlength\textfloatsep{5pt}
\begin{figure}[t]
\centering
\includegraphics[width=\columnwidth]{icra2018_attractor_v2.eps}
\caption{Embedding attractor structure into the internal state of RNN}
\label{fig_kase_attractor}
\hspace{10mm}
\end{figure}



\subsection{Integrated representation of language and motion}
Finally, we describe the integrated language and motion learning method for complex motion generation. The correspondence between language, which is a discrete representation, and robot motion in the continuous space implies changes depending on the past context and current situation. For example, even for a simple instruction such as ``hit the blue bell,'' the motion to accomplish this differs depending on the location of the blue bell or the robot's posture while waiting. The robot needs to acquire an internal representation in a form that can realize such a polysemous and unique grounding each time, i.e., a form that can be generalized even in unknown situations.

Deep predictive learning embeds the understanding of language instructions and motion generation into a set of attractor structures. In \cite{Yamada2016fn}, an RNN predictively learns a sequence of multiple iterations of episodes comprising waiting, language instruction, and robot motion generation, thereby acquiring a circumferential attractor representing the above episodic iterations as the internal dynamics of the model (Figure~\ref{fig_language_reduce}). When a language command is input to the robot waiting at the fixed point attractor, it is integrated with the environmental information given as visual information, and the internal state of the RNN transitions in a branching manner according to the situation. The robot's motion is then generated according to the meaning of the instruction in the environment. Finally, the internal state reconverges to the fixed point and waits for the next instruction. The generalization ability to generate appropriate motions for new situations without teaching all combinations of environmental states and instructions is demonstrated.

%-------------------- Fig --------------------
\setlength\textfloatsep{5pt}
\begin{figure}[t]
\centering
\includegraphics[width=\columnwidth]{rsj2022_language_reduce.eps}
\caption{Grounding of language instructions and motion generation in the internal dynamics of RNNs}
\label{fig_language_reduce}
\hspace{10mm}
\end{figure}

The context that can be learned as a precondition for motion generation includes environmental information and the robot state. In \cite{Yamada2017fn}, a flag-raising game is learned using the framework in \cite{Yamada2016fn}. Fixed point attractors, corresponding to the four standby postures of the robot (up-up, up-down, down-up, down-down), are structured as the internal dynamics of the RNN. When instructions are given, the internal state is raised according to the flag color at that time and converges to one of the fixed points while generating appropriate motions. It is expected to be deployed in use cases where robots work in response to sequential instructions from humans.

Furthermore, the model can be extended to a format that can be used for bidirectional conversion, i.e., conversion from motion to verbal description and vice versa. In \cite{Yamada2018ral}, two autoencoders, for language and sensorimotor, were prepared and trained with the constraint that the feature vectors are close in the corresponding sentences and movements. After training, the model realized bi-directional conversion from linguistic instructions to robot motion and from robot motion to linguistic explanations via feature vectors. Furthermore, \cite{Toyoda2021ral} used a distributed representation conversion module for words using the Word2vec model, which has recently attracted attention in natural language processing, to generate appropriate motions for instructions containing novel words similar to the trained words.


In the example above, we introduced a method for integrating robot motions by giving instructions as language. The language instructions are given in word units, but a generic planning description language (e.g., planning domain description language (PDDL)) used in conventional automatic planning methods can also be used. In \cite{Kase2020icra}, the method is learned by grounding images with descriptive expressions that define the environment using the description language (Figure~\ref{fig_kase2022icra}). Based on the discrete symbolic representations extracted from the images, designing motions using arbitrary automatic planning methods is possible. Conventional methods, such as PDDL, plan in advance, making it difficult to respond to environmental changes, but dynamic planning methods allow to dynamically determine the robot's motion. In addition, this research aims to realize a robust system using a deep learner to generate motions based on the motion planning guidelines.

%-------------------- Fig --------------------
\setlength\textfloatsep{5pt}
\begin{figure}[t]
\centering
\includegraphics[width=\columnwidth]{icra2020_overview.eps}
\caption{Grounding with generic planning description language and feature extraction}
\label{fig_kase2022icra}
\hspace{10mm}
\end{figure}    


\subsection{Examples of various task}
We have introduced several methods for complex motion generation. We have applied deep predictive learning, which is the basis of these methods, to various robot or object manipulation tasks to verify the generality of the concept. One difficult task with conventional robot technology is manipulating irregularly shaped objects, such as cloth or string, whose shape changes upon contact, making it difficult to describe the entire situation programmatically. To address this problem, \cite{Koma2016ral} generalized the robot's experience by training the robot to fold several towels and performing integrated learning based on the acquired image-motion information. Moreover, from the limited training data, the robot repeatedly performed motions corresponding to object misalignment and different objects (such as books). Unlike conventional methods for handling flexible objects, the robot could generate movements naturally without modeling, image processing, or other costly processing\footnote{https://youtu.be/YH1TrL1q6Po}.

We have since extended the deep predictive learning model to realize a knot--tying ~\cite{Suzuki2021iros} and button-hanging ~\cite{Fujii2022sii} tasks with the same model. We confirmed that including detailed object information such as proximity and tactile sensors as model inputs improved prediction performance ~\cite{Suzuki2021iros,Ito2022sii}. In particular, \cite{ichiwara2022contact} realized a task of unzipping a cloth bag of various shapes and colors by introducing an attention mechanism that extracts useful information (such as the position and direction of the object) from the robot's visual and tactile information. Because the attention mechanism focuses only on important details during task execution, it can generate robust movements in response to environmental and lighting conditions.

In addition, we are developing a similar hierarchical prediction model for use in robotic simulations of neurodevelopmental disorders.
Specifically, by reproducing functional disconnections between hierarchical layers \cite{idei2021paradoxical} and excitatory-inhibitory imbalances \cite{idei2020homogeneous} of neurons, it has been shown that real robots can produce abnormalities in sensory accuracy and the prediction accuracy of upper layers, observed as behavioral abnormalities.




\section{Conclusion}

This paper introduces the background technology of our proposed deep predictive learning, such as predictive coding, related research, and robot applications. We also describe our approach and application examples for complex robot motion generation using deep predictive learning from the following three perspectives.


\begin{enumerate}
\renewcommand{\labelenumi}{\roman{enumi}}
\item Switching multiple prediction models based on prediction error.
\item Combining multiple motions by designing the dynamics inside the model.
\item Acquiring integrated representation of language-motion and its application to complex motion generation.
\end{enumerate}


Although each method has been developed separately, they share many common elements and are expected to be integrated in the future. However, some issues need to be addressed, such as training data collection and segmentation methods, to improve accuracy and apply the methods to challenging tasks. In addition, deep predictive learning highly depends on the learned physical characteristics. This can be addressed by transfer learning or other methods when there is some similarity in physical characteristics, but it cannot be applied to robots with completely different structures. We plan to revisit the concept of deep predictive learning from the perspective of cognitive robotics, where the concept originated, and apply the method to a longer-term and more diverse set of tasks better suited to real-world problems.


\section*{Acknowledgement}
We thank Hayato Idei for his helpful comments on the description of predictive coding. This research was supported by JST Moonshot-type R\&D Project (JPMJMS2031), JST ACT-X (JPMJAX190I), and Hitachi, Ltd. We would like to express our gratitude to them.

\newpage
\bibliographystyle{IEEEtran}
\bibliography{ref}

\end{document}